\chapter{Max-stable vectors and processes}
\label{chap:max-stable}

We follow de Haan, \emph{A spectral representation for max-stable processes}, Ann.\ Probab.\ 12
(1984) [dH84], [dHF, Ch.~6 and~9] and [Resnick, Ch.~5].

\begin{definition}[Simple max-stable vector]
  \label{def:simple-max-stable}
  \uses{def:max-stable}
  \lean{ExtremeValueTheory.IsSimpleMaxStable, ExtremeValueTheory.maxStableExponent}
  \leanok
  A law $\mu$ on $\bR^d$ is simple max-stable if it has unit Fr\'echet margins
  $\PR[Y_i \le y] = e^{-1/y}$ and $\PR[Y \le y]^n = \PR[Y \le y/n]$ for $n \ge 1$, $y > 0$. Its exponent
  function is $V(y) = -\log \PR[Y \le y]$ [dHF, \S6.1].
\end{definition}

\begin{theorem}[Homogeneity of the exponent function]
  \label{thm:max-stable-exponent}
  \uses{def:simple-max-stable}
  \lean{ExtremeValueTheory.IsSimpleMaxStable.measure_Iic_pos, ExtremeValueTheory.IsSimpleMaxStable.measureReal_Iic_smul,
    ExtremeValueTheory.IsSimpleMaxStable.maxStableExponent_smul, ExtremeValueTheory.IsSimpleMaxStable.exp_neg_maxStableExponent,
    ExtremeValueTheory.IsSimpleMaxStable.inv_le_maxStableExponent, ExtremeValueTheory.IsSimpleMaxStable.maxStableExponent_le_sum_inv}
  \leanok
  $\PR[Y \le ty] = \PR[Y \le y]^{1/t}$ for all real $t > 0$, so $V(ty) = V(y)/t$, and
  $\max_i 1/y_i \le V(y) \le \sum_i 1/y_i$.
\end{theorem}
\begin{proof}
  \leanok
  Rational $t$ from max-stability, then monotonicity and continuity.
\end{proof}

\begin{theorem}[Spectral construction of max-stable vectors]
  \label{thm:max-stable-spectral}
  \uses{def:simple-max-stable, thm:prm-exists, thm:prm-uniqueness}
  \lean{ExtremeValueTheory.spectralMax, ExtremeValueTheory.spectralVector, ProbabilityTheory.IsPoissonRandomMeasure.measure_apply_eq_zero, ProbabilityTheory.IsPoissonRandomMeasure.measure_forall_essSup_le,
    ExtremeValueTheory.radialMeasure_prod_iUnion, ProbabilityTheory.IsPoissonRandomMeasure.measure_spectralVector_le,
    ProbabilityTheory.IsPoissonRandomMeasure.isSimpleMaxStable_map_spectralVector, ExtremeValueTheory.exists_isSimpleMaxStable_of_spectral,
    ExtremeValueTheory.exists_isSimpleMaxStable_of_spectralMeasure}
  \leanok
  Let $N$ be a Poisson random measure on $(0, \infty) \times S$ with intensity $\zeta^{-2} \ud\zeta \otimes \mu$ and
  $f_i \ge 0$ with $\int f_i \ud\mu = 1$. Then $Y_i = \max_{(\zeta, s) \in N} \zeta f_i(s)$ is simple max-stable with
  $\PR[Y \le y] = \exp(-\int \max_i f_i(s)/y_i \, \mu(\ud s))$. In particular every measure $H$ with
  $\int w_i \ud H = 1$ is the spectral measure of a simple max-stable law [dH84], [Resnick,
  Proposition~5.11]. The converse (every simple max-stable law has a spectral measure) is not
  formalized.
\end{theorem}
\begin{proof}
  \leanok
  $\{Y \le y\}$ is the event that $N$ has no point in $\bigcup_i \{\zeta f_i(s) > y_i\}$, whose intensity is
  $\int \max_i f_i/y_i \ud\mu$.
\end{proof}

\begin{theorem}[Spectral representation of max-stable processes]
  \label{thm:max-stable-process}
  \uses{thm:max-stable-spectral}
  \lean{ProbabilityTheory.IsPoissonRandomMeasure.isSimpleMaxStable_map_process, ProbabilityTheory.IsPoissonRandomMeasure.measure_process_le, ExtremeValueTheory.smithKernel,
    ProbabilityTheory.IsPoissonRandomMeasure.isSimpleMaxStable_smith, ExtremeValueTheory.independentSpectralMeasure, ExtremeValueTheory.comonotoneSpectralMeasure,
    ProbabilityTheory.IsPoissonRandomMeasure.measure_spectralVector_le_independent, ProbabilityTheory.IsPoissonRandomMeasure.measure_spectralVector_le_comonotone}
  \leanok
  For $W \colon T \to S \to [0, \infty]$ with $\int W(t) \ud\mu = 1$, the process $Y(t) = \max_{(\zeta,s) \in N} \zeta W(t)(s)$
  has simple max-stable finite-dimensional laws with
  $\PR[Y(t_i) \le y_i \ \forall i] = \exp(-\int \max_i W(t_i)/y_i \ud\mu)$ [dH84]. Examples: the Smith model
  with a Gaussian kernel; independence (spectral mass on the vertices) gives
  $\prod_i e^{-1/y_i}$, complete dependence gives $\min_i e^{-1/y_i}$.
\end{theorem}
\begin{proof}
  \leanok
  Apply Theorem~\ref{thm:max-stable-spectral} to the finitely many functions $W(t_i)$.
\end{proof}

\begin{theorem}[Every simple max-stable law has a spectral measure]
  \label{thm:max-stable-converse}
  \uses{thm:max-stable-spectral, thm:max-stable-exponent}
  \lean{ExtremeValueTheory.IsSimpleMaxStable.exists_hasSpectralMeasure, ExtremeValueTheory.HasSpectralMeasure.measure_univ_toReal,
    ExtremeValueTheory.IsSimpleMaxStable.ae_forall_pos, ExtremeValueTheory.IsSimpleMaxStable.convexOn_pickandsDependence,
    ExtremeValueTheory.IsSimpleMaxStable.exists_pickandsDependence_eq_pickandsSpectral}
  \leanok
  For every finite dimension, a simple max-stable law has a finite spectral measure $H$ on the unit sphere of the sup
  norm with $\int w_i \ud H = 1$ and $\PR[Y \le y] = \exp(-\int \max_i w_i/y_i \ud H)$ [dHF, Theorem~6.1.9],
  [Resnick, Proposition~5.11]; its total mass is the extremal coefficient. Consequently the Pickands dependence
  function of every bivariate simple max-stable law is convex.
\end{theorem}
\begin{proof}
  \leanok
  With $R = \max_i Y_i$ and $\Theta = Y/R$, homogeneity gives $\PR[\max_i Y_i/y_i > r] = 1 - e^{-V(y)/r}$. The finite measures
  $\Lambda_M = \E[\min(R, M);\, \Theta \in \cdot]/\log M$ on the compact cube satisfy $\int \max_i \theta_i/y_i \ud\Lambda_M \to V(y)$; a weak
  cluster point (Prokhorov) is the spectral measure.
\end{proof}

\begin{theorem}[Uniqueness of the spectral measure]
  \label{thm:max-stable-unique}
  \uses{thm:max-stable-converse}
  \lean{ExtremeValueTheory.eq_of_lintegral_iSup_div_eq, ExtremeValueTheory.spectralExponentMeasure, ExtremeValueTheory.lintegral_iSup_div_eq_iInf,
    ExtremeValueTheory.HasSpectralMeasure.unique, ExtremeValueTheory.IsSimpleMaxStable.existsUnique_hasSpectralMeasure}
  \leanok
  A finite measure on the sup-norm sphere is determined by the integrals $\int \max_i w_i/y_i \ud H$, $y > 0$; hence every simple
  max-stable law has exactly one spectral measure on the sphere [dHF, \S6.1].
\end{theorem}
\begin{proof}
  \leanok
  The exponent measure $\nu_H$ (image of $r^{-2}\ud r \otimes H$) is determined on complements of lower orthants, hence by
  inclusion--exclusion on the $\pi$-system of upper orthants, and $H$ is its angular part.
\end{proof}

\begin{theorem}[Simplex form and the angular measure]
  \label{thm:max-stable-simplex}
  \uses{thm:max-stable-unique, def:pickands-dependence}
  \lean{ExtremeValueTheory.normTransport, ExtremeValueTheory.simplexSpectralMeasure, ExtremeValueTheory.HasSpectralMeasure.simplex, ExtremeValueTheory.simplexSpectralProb,
    ExtremeValueTheory.eq_of_lintegral_iSup_div_eq_of_simplex, ExtremeValueTheory.pickandsSpectral_simplexSpectralMeasure,
    ExtremeValueTheory.pickandsAngularMeasure, ExtremeValueTheory.HasSpectralMeasure.pickandsDependence_eq_angular,
    ExtremeValueTheory.HasSpectralMeasure.pickandsDependence_one_sub_eq_angular}
  \leanok
  Transporting $H$ to the sum-norm simplex gives the standard spectral measure (total mass $d$, $\int v_i = 1$), also unique. In
  two dimensions, with the angular measure $Q$ on $[0,1]$ (mass $2$, mean $1$),
  $A(1-t) = \int \max((1-t)q, t(1-q)) \ud Q(q)$, the convention of Beirlant et al.\ (8.48).
\end{theorem}
\begin{proof}
  \leanok
  The change of norm $w \mapsto w/N(w)$ with density $N(w)$ preserves integrals of positively homogeneous functions.
\end{proof}

\begin{theorem}[Max-stable vectors are regularly varying]
  \label{thm:max-stable-rv}
  \uses{thm:max-stable-unique, def:rv-vector}
  \lean{ExtremeValueTheory.maxStableTailMeasure, ExtremeValueTheory.IsSimpleMaxStable.isRegVarVector, ExtremeValueTheory.spectralMeasure_maxStableTailMeasure,
    ExtremeValueTheory.IsSimpleMaxStable.exists_isRegVarVector}
  \leanok
  A simple max-stable vector is regularly varying with index $1$; its exponent measure is $\nu_H$ normalized, and its spectral
  measure in the sense of multivariate regular variation is $H/H(\text{univ})$, where $H(\text{univ})$ is the extremal
  coefficient.
\end{theorem}
\begin{proof}
  \leanok
  Exact tails $\PR[\exists i,\ X_i/x > d_i] = 1 - e^{-\Phi(d)/x}$; inclusion--exclusion on boxes, rational boxes for open sets, and
  the portmanteau theorem.
\end{proof}
